% from ega1-10-fr.tex
\oldpage[I]{187}
\emph{d'Idéaux de définition} s'il existe un recouvrement
$(U_\alpha)$ de $\mathfrak{X}$ par des ouverts formels affines tel que,
pour tout $\alpha$, la famille des
$\mathscr{J}_\lambda|U_\alpha$ soit un système fondamental d'Idéaux de
définition (\hyperref[I.10.3.6-fr]{10.3.6}) du schéma formel affine
induit par $\mathfrak{X}$ sur $U_\alpha$. Il résulte de la remarque finale
de (\hyperref[I.10.3.7-fr]{10.3.7}) que lorsque $\mathfrak{X}$ est un
schéma formel affine, cette définition coïncide avec la définition donnée
dans (\hyperref[I.10.3.7-fr]{10.3.7}). Pour tout ouvert $V$ de
$\mathfrak{X}$, les restrictions $\mathscr{J}_\lambda|V$ forment alors un
système fondamental d'Idéaux de définition du préschéma formel induit sur
$V$, en vertu de (\hyperref[I.10.3.1.1-fr]{10.3.1.1}). Si
$\mathfrak{X}$ est un préschéma formel \emph{localement noethérien}, et
$\mathscr{J}$ un Idéal de définition de $\mathfrak{X}$, il résulte de
(\hyperref[I.10.3.6-fr]{10.3.6}) que les puissances $\mathscr{J}^n$
forment un système fondamental d'Idéaux de définition de $\mathfrak{X}$.
\end{env}

\begin{env}[10.5.2]
\label{I.10.5.2-fr}
Soient $\mathfrak{X}$ un préschéma formel, $\mathscr{J}$ un Idéal de
définition de $\mathfrak{X}$. Alors l'espace annelé
$(\mathfrak{X},\mathscr{O}_{\mathfrak{X}}/\mathscr{J})$ est un
\emph{préschéma} (usuel), qui est affine (resp. localement noethérien,
resp. noethérien) lorsque $\mathfrak{X}$ est un schéma formel affine
(resp. un préschéma formel localement noethérien, resp. noethérien)~; on
est en effet aussitôt ramené au cas affine, et alors la proposition a déjà
été démontrée dans (\hyperref[I.10.3.2-fr]{10.3.2}). En outre, si
$\theta:\mathscr{O}_{\mathfrak{X}}\to
\mathscr{O}_{\mathfrak{X}}/\mathscr{J}$ est l'homomorphisme canonique,
$u=(1_{\mathfrak{X}},\theta)$ est un \emph{morphisme} (dit
\emph{canonique}) de préschémas formels
$(\mathfrak{X},\mathscr{O}_{\mathfrak{X}}/\mathscr{J})\to
(\mathfrak{X},\mathscr{O}_{\mathfrak{X}})$, car ici encore, cela a été vu
dans le cas affine (\hyperref[I.10.3.2-fr]{10.3.2}), auquel on se ramène
aussitôt.
\end{env}

\begin{proposition}[10.5.3]
\label{I.10.5.3-fr}
Soient $\mathfrak{X}$ un préschéma formel, $(\mathscr{J}_\lambda)$ un
système fondamental d'Idéaux de définition de $\mathfrak{X}$. Alors le
faisceau d'anneaux topologiques $\mathscr{O}_{\mathfrak{X}}$ est limite
projective des faisceaux d'anneaux pseudo-discrets
(\hyperref[0.3.8.1-fr]{0, 3.8.1})
$\mathscr{O}_{\mathfrak{X}}/\mathscr{J}_\lambda$.
\end{proposition}

\begin{proof}
Comme la topologie de $\mathfrak{X}$ admet une base d'ouverts formels
affines quasi-compacts (\hyperref[I.10.4.3-fr]{10.4.3}), on est ramené au
cas affine, où la proposition est conséquence de
(\hyperref[I.10.3.5-fr]{10.3.5}),
(\hyperref[I.10.3.2-fr]{10.3.2}) et de la définition
(\hyperref[I.10.1.1-fr]{10.1.1}).
\end{proof}

Il n'est pas certain que tout préschéma formel admette des Idéaux de
définition. Toutefois~:

\begin{proposition}[10.5.4]
\label{I.10.5.4-fr}
Soit $\mathfrak{X}$ un préschéma formel localement noethérien. Il existe
un plus grand Idéal de définition $\mathscr{T}$ de $\mathfrak{X}$~;
c'est le seul Idéal de définition $\mathscr{J}$ tel que le préschéma
$(\mathfrak{X},\mathscr{O}_{\mathfrak{X}}/\mathscr{J})$ soit réduit. Si
$\mathscr{J}$ est un Idéal de définition de $\mathfrak{X}$,
$\mathscr{T}$ est l'image réciproque par
$\mathscr{O}_{\mathfrak{X}}\to
\mathscr{O}_{\mathfrak{X}}/\mathscr{J}$ du Nilradical de
$\mathscr{O}_{\mathfrak{X}}/\mathscr{J}$.
\end{proposition}

\begin{proof}
Supposons d'abord que $\mathfrak{X}=\operatorname{Spf}(A)$, où $A$ est
un anneau adique noethérien. L'existence et les propriétés de
$\mathscr{T}$ résultent immédiatement de
(\hyperref[I.10.3.4-fr]{10.3.4}) et
(\hyperref[I.5.1.1-fr]{5.1.1}), compte tenu de l'existence et des
propriétés du plus grand idéal de définition de $A$
(\hyperref[0.7.1.6-fr]{0, 7.1.6} et
\hyperref[0.7.1.7-fr]{7.1.7}).

Pour prouver l'existence et les propriétés de $\mathscr{T}$ dans le cas
général, il suffit de montrer que si $U\supset V$ sont deux ouverts
formels affines noethériens de $\mathfrak{X}$, le plus grand Idéal de
définition $\mathscr{T}_U$ de $U$ induit le plus grand Idéal de définition
$\mathscr{T}_V$ de $V$~; mais comme
$(V,(\mathscr{O}_{\mathfrak{X}}|V)/(\mathscr{T}_U|V))$ est réduit, cela
résulte de ce qui précède.
\end{proof}

On désigne par $\mathfrak{X}_{\mathrm{red}}$ le préschéma (usuel) réduit
$(\mathfrak{X},\mathscr{O}_{\mathfrak{X}}/\mathscr{T})$.

\begin{corollary}[10.5.5]
\label{I.10.5.5-fr}
Soient $\mathfrak{X}$ un préschéma formel localement noethérien,
$\mathscr{T}$ le plus grand Idéal de définition de $\mathfrak{X}$~; pour
tout ouvert $V$ de $\mathfrak{X}$, $\mathscr{T}|V$ est le plus grand
Idéal de définition du préschéma formel induit par $\mathfrak{X}$ sur
$V$.
\end{corollary}

